% Level 3 workbook, Week 8 -- Data Detectives
\chapterweek{8}{Data Detectives}{Data tells stories when you clean it, count it and chart it.}

\section{What you will learn}
\begin{itemize}[itemsep=1pt]
  \item How to read a \textbf{CSV file} with Python's \texttt{csv} module.
  \item How to \textbf{clean} messy data: extra spaces, MiXeD capitals, empty rows, missing numbers.
  \item How to \textbf{count} with a dictionary and work out an \textbf{average}.
  \item How to draw a \textbf{bar chart} that everyone can read -- and how to keep survey answers private.
\end{itemize}

\section{The big idea: the data detective's pipeline}
A data detective always works in the same order. Skipping a step gives wrong answers!

\begin{center}
%% alt: A five-step data pipeline. 1 Read the CSV file into rows. 2 Clean: spaces, capitals and empty rows. 3 Count with a dictionary tally. 4 Calculate averages. 5 Chart with bars and labels.
\begin{tikzpicture}[stage/.style={draw,rounded corners=3pt,minimum height=1.1cm,text width=2.3cm,align=center,font=\small},>=Stealth,node distance=4mm]
  \node[stage,fill=kcsky!25] (a) {\textbf{1 Read}\\CSV file $\rightarrow$ rows};
  \node[stage,fill=kcorange!25,right=of a] (b) {\textbf{2 Clean}\\spaces, capitals, empty rows};
  \node[stage,fill=kcgreen!20,right=of b] (c) {\textbf{3 Count}\\dictionary tally};
  \node[stage,fill=kcyellow!40,right=of c] (d) {\textbf{4 Calculate}\\averages};
  \node[stage,fill=kcpurple!25,right=of d] (e) {\textbf{5 Chart}\\bars + labels};
  \draw[->,thick] (a) -- (b); \draw[->,thick] (b) -- (c); \draw[->,thick] (c) -- (d); \draw[->,thick] (d) -- (e);
\end{tikzpicture}
\end{center}

\subsection{1. Read: what is a CSV file?}
\textbf{CSV} means \emph{comma-separated values}: a table saved as plain text. The first line is the \textbf{header} (the column names). Every other line is one \textbf{row}. Here is the start of \texttt{favorite\_games.csv} -- made-up, anonymous survey answers (a response number, no names):

\begin{center}\small
\begin{tabular}{@{}l@{}}
\texttt{response,age\_group,favorite\_game,genre,hours\_per\_week,stars}\\
\texttt{1,10-11,Block Craft,building,5,5}\\
\texttt{2,12-13, block craft ,building,3,4}\\
\texttt{5,,,,,}\\
\texttt{7,12-13,Kart Dash,racing,,3}\\
\end{tabular}
\end{center}

\py|csv.DictReader| turns each row into a \textbf{dictionary}: \py|row["favorite_game"]| is \py|"Block Craft"|.
\begin{pycode}
import csv
with open("favorite_games.csv", newline="") as f:
    for row in csv.DictReader(f):
        print(row["favorite_game"], row["stars"])
\end{pycode}

\subsection{2. Clean: wash the data first}
Every value arrives as \textbf{text}, even numbers: \py|"5"|, not \py|5|. And people type things in different ways. Spot the mess:

\begin{center}\small
\begin{tabular}{@{}p{3.3cm}p{4.1cm}p{\dimexpr\linewidth-7.4cm-4\tabcolsep}@{}}
\toprule
\textbf{Messy value} & \textbf{Problem} & \textbf{Clean it with}\\ \midrule
\texttt{" block craft "} & extra spaces, small letters & \py|.strip().title()| $\rightarrow$ \texttt{"Block Craft"}\\
\texttt{"BLOCK CRAFT"} & capitals & \py|.title()| $\rightarrow$ \texttt{"Block Craft"}\\
\texttt{"5"} & a number stored as text & \py|int("5")| $\rightarrow$ \texttt{5}\\
\texttt{""} (hours) & missing answer & use \py|None|, and skip it in averages\\
\texttt{5,,,,,} & empty row & skip the whole row\\
\bottomrule
\end{tabular}
\end{center}

\subsection{3 and 4. Count and calculate}
Counting with a dictionary is like a tally chart. The first time you see a game, \py|counts.get(game, 0)| gives 0:
\begin{pycode}
counts = {}
for row in rows:
    game = row["favorite_game"]
    counts[game] = counts.get(game, 0) + 1
# {"Block Craft": 5, "Kart Dash": 4, ...}
\end{pycode}
The \textbf{average} (mean) is the total divided by how many: \py|sum(numbers) / len(numbers)|. Leave out missing answers -- counting a blank as 0 would make the average unfairly small.

\subsection{5. Chart: readable by everyone}
A good chart never relies on color alone. Write the \textbf{name} under each bar and the \textbf{number} on top, so someone who sees colors differently can still read it.
\begin{center}
%% alt: A bar chart of favorite games: Block Craft 5, Kart Dash 4 and Puzzle Drop 3.
\begin{tikzpicture}[x=2.2cm,y=0.45cm]
  \foreach \x/\h/\name/\c in {0/5/Block Craft/kcorange,1/4/Kart Dash/kcsky,2/3/Puzzle Drop/kcgreen}{
    \draw[fill=\c!60] (\x,0) rectangle (\x+0.7,\h);
    \node[above,font=\small] at (\x+0.35,\h) {\h};
    \node[below,font=\footnotesize] at (\x+0.35,0) {\name};
  }
  \draw (-0.3,0) -- (2.9,0);
\end{tikzpicture}
\end{center}

\begin{exercise}[Survey cleaner]
Folder: \texttt{week08/ex-survey-cleaner}. Open \texttt{favorite\_games.csv} in Thonny and circle the mess you find. Then finish \texttt{to\_number()}, \texttt{is\_empty()} and add \py|.title()| to the game names. \textbf{Done:} \texttt{Clean answers: 17}.

How many rows did the cleaner throw away, and why? \answerlines{1}
\end{exercise}

\begin{exercise}[Survey stats]
Folder: \texttt{week08/ex-survey-stats}. Finish \texttt{count\_by()}, \texttt{average()} and \texttt{most\_popular()}. Then answer from your program's report:

Most popular game: \rule{4cm}{0.4pt} \quad Average hours per week: \rule{2cm}{0.4pt}

Adventure games got 5.0 stars on average from only 2 answers. Should we trust that? Why? \answerlines{1}
\end{exercise}

\begin{exercise}[Bar chart]
Folder: \texttt{week08/ex-bar-chart}. Make \texttt{bar\_height()} scale the bars so the biggest is \texttt{MAX\_HEIGHT} pixels, and write a number on top and a name under every bar. \textbf{Done:} six labeled bars from Block Craft (5) to Space Chess (1).
\end{exercise}

\begin{safety}[Consent and privacy in surveys]
Before collecting data, ask: \emph{Did people agree to answer? Can they skip a question? Do I really need this question?} Never collect real names, addresses, schools, birthdays or photos. Even without names, answers together can point to one person (\enquote{the only 13-year-old in the chess club}). Before sharing data, \textbf{anonymize} it: keep only the columns you need and give each row an ID like R1, R2.
\end{safety}

\begin{vocab}
  \vocabword{CSV}{Comma-separated values: a table saved as plain text, one row per line.}
  \vocabword{header}{The first line of a CSV file, with the column names.}
  \vocabword{data cleaning}{Fixing messy data (spaces, capitals, missing or empty values) before using it.}
  \vocabword{average (mean)}{The total divided by how many values there are.}
  \vocabword{sample}{The group of people who answered. A small sample can give misleading results.}
  \vocabword{consent}{Agreeing, freely, to take part -- with the choice to say no.}
  \vocabword{anonymize}{Remove or replace anything that could show who gave an answer.}
\end{vocab}

\begin{reflection}
What story does the bar chart tell? What can it \emph{not} tell us?
\answerlines{2}
Would you answer a survey that asked for your full name? Why or why not?
\answerlines{1}
\end{reflection}
